% part: intuitionistic-logic
% chapter: introduction
% section: bhk-interpretation

\documentclass[../../../include/open-logic-section]{subfiles}

\begin{document}

\olfileid[hi]{int}{int}{bhk}

\olsection{ब्राउवर--हाइटिंग--कोल्मोगोरोव व्याख्या}

\begin{editorial}
  BHK व्याख्या का उपयोग करके अंतःप्रज्ञावादी प्रतिज्ञप्तियों की मान्यता के
  प्रमाण उलझाऊ हैं; उन्हें अधिक स्पष्ट ढंग से समझाना आवश्यक है।
\end{editorial}

अंतःप्रज्ञावादी संयोजकों की एक अनौपचारिक रचनात्मक व्याख्या है, जिसे सामान्यतः
ब्राउवर--हाइटिंग--कोल्मोगोरोव व्याख्या कहते हैं। इसमें ``रचना'' की धारणा का
उपयोग होता है, जिसे आप रचनात्मक प्रमाण समझ सकते हैं। (BHK व्याख्या में हम
``प्रमाण'' शब्द का उपयोग नहीं करते, ताकि उसे किसी औपचारिक
!!{derivation}-प्रणाली की !!{derivation} की धारणा से भ्रमित न किया जाए।)
इस सहज धारणा के आधार पर BHK व्याख्या अंतःप्रज्ञावादी संयोजकों के अर्थ समझाती
है।

\begin{enumerate}
\item हम मानते हैं कि हमें ज्ञात है कि परमाण्विक कथन की रचना किसे माना जाता
  है।
\item $!A_1 \land !A_2$ की रचना ऐसा युग्म $\tuple{M_1, M_2}$ है जिसमें
  $M_1$,~$!A_1$ की और $M_2$,~$!A_2$ की रचना है।
\item $!A_1 \lor !A_2$ की रचना ऐसा युग्म $\tuple{s, M}$ है जिसमें या तो
  $s$ का मान~$1$ और $M$,~$!A_1$ की रचना है, अथवा $s$ का मान~$2$ और
  $M$,~$!A_2$ की रचना है।
\item $!A \lif !B$ की रचना ऐसा फलन है जो~$!A$ की किसी रचना को~$!B$ की रचना
  में बदल देता है।
\item $\lfalse$ (असंगति) की कोई रचना नहीं है।
\item $\lnot !A$ को $!A \lif \lfalse$ का पर्याय परिभाषित किया जाता है।
  अर्थात $\lnot !A$ की रचना ऐसा फलन है जो~$!A$ की रचना को~$\lfalse$ की
  रचना में बदल देता है।
\end{enumerate}

\begin{ex}
उदाहरण के लिए $\lnot \lfalse$ लें। उसकी रचना ऐसा फलन है जो आगत के रूप में
$\lfalse$ की कोई भी रचना दिए जाने पर निर्गत में~$\lfalse$ की रचना देता है।
स्पष्टतः तत्समक फलन~$\Id{}$ ऐसी रचना है: $\lfalse$ की कोई रचना~$M$ दिए
जाने पर $\Id{}(M) = M$,~$\lfalse$ की रचना देता है।
\end{ex}

सामान्यतः $\lnot !A$ का अर्थ है, ``$!A$ की रचना असंभव है।''

\begin{ex}
किसी भी प्रतिज्ञप्ति~$!A$ के लिए $!A \lif \lnot \lnot !A$, अर्थात
$!A \lif ((!A \lif \lfalse) \lif \lfalse)$, सिद्ध करें। इसकी रचना ऐसा
फलन~$f$ होना चाहिए जो $!A$ की कोई रचना~$M$ दिए जाने पर
$(!A \lif \lfalse) \lif \lfalse$ की रचना~$f(M)$ लौटाए। $f$ इस रचना को
इस प्रकार बनाता है: हमें ऐसा फलन~$g$ परिभाषित करना है जो आगत में
$!A \lif \lfalse$ की कोई रचना~$h$ मिलने पर निर्गत में~$\lfalse$ की रचना
दे। हम $g$ को इस प्रकार परिभाषित कर सकते हैं: आगत~$h$ को $!A$ की उस
रचना~$M$ पर लागू करें जो हमें पहले मिली थी। चूँकि $h$ का निर्गत~$h(M)$,
$\lfalse$ की रचना है, इसलिए यदि $M$,~$!A$ की रचना है तो
$f(M)(h) = h(M)$,~$\lfalse$ की रचना है।
\end{ex}

\begin{ex}
$\lnot(!A \land \lnot !A)$, अर्थात
$(!A \land (!A \lif \lfalse)) \lif \lfalse$, की रचना दें। यह ऐसा फलन~$f$
है जो आगत में $!A \land (!A \lif \lfalse)$ की कोई रचना~$M$ मिलने पर
$\lfalse$ की रचना देता है। किसी संयोजन $!B_1 \land !B_2$ की रचना ऐसा युग्म
$\tuple{N_1, N_2}$ है जिसमें $N_1$,~$!B_1$ की और $N_2$,~$!B_2$ की रचना
है। हम ऐसे फलन $p_1$ और $p_2$ परिभाषित कर सकते हैं जो
$!B_1 \land !B_2$ की किसी रचना से क्रमशः $!B_1$ और $!B_2$ की रचनाएँ
प्राप्त करते हैं:
\begin{align*}
  p_1(\tuple{N_1, N_2}) &= N_1\\
  p_2(\tuple{N_1, N_2}) &= N_2
\end{align*}
$f$ इस प्रकार कार्य करता है: पहले वह अपने आगत~$M$ पर $p_1$ लगाता है। इससे
$!A$ की रचना मिलती है। फिर वह~$M$ पर $p_2$ लगाता है, जिससे
$!A \lif \lfalse$ की रचना मिलती है। ऐसी रचना स्वयं एक फलन~$p_2(M)$ है,
जो आगत में~$!A$ की कोई रचना मिलने पर~$\lfalse$ की रचना देता है। दूसरे
शब्दों में, $p_2(M)$ को $p_1(M)$ पर लगाने से हमें~$\lfalse$ की रचना मिलती
है। अतः हम $f(M) = p_2(M)(p_1(M))$ परिभाषित कर सकते हैं।
\end{ex}

\begin{ex}
$((!A \land !B) \lif !C) \lif (!A \lif (!B \lif !C))$ की रचना दें;
अर्थात ऐसा फलन~$f$ जो $(!A \land !B) \lif !C$ की किसी रचना~$g$ को
$(!A \lif (!B \lif !C))$ की रचना में बदल दे। रचना~$g$ स्वयं एक फलन है
($!A \land !B$ की रचनाओं से~$C$ की रचनाओं तक)। और निर्गत~$f(g)$ ऐसा
फलन~$h_g$ है जो $!A$ की रचनाओं को उन फलनों में भेजता है जो~$!B$ की रचनाओं
को~$!C$ की रचनाओं में भेजते हैं।

ठीक है, यह उलझाऊ है। हमें कोई निश्चित फलन~$h_g$ बनाना है, जो आगत~$g$ के
लिए $f$ का निर्गत होगा। $h_g$ का आगत~$!A$ की कोई रचना~$M$ है। $h_g(M)$
का निर्गत ऐसा फलन~$k_M$ होना चाहिए जो~$!B$ की रचनाओं~$N$ को~$!C$ की रचनाओं
में भेजता है। $k_{g,M}(N) = g(\tuple{M, N})$ लें। स्मरण रखें कि
$\tuple{M, N}$,~$!A \land !B$ की रचना है। अतः $k_{g,M}$,
$!B \lif !C$ की रचना है: वह~$!B$ की रचनाओं~$N$ को~$!C$ की रचनाओं में भेजता
है। अब $h_g(M) = k_{g,M}$ लें। यह ऐसा फलन है जो~$!A$ की रचनाओं~$M$ को
$!B \lif !C$ की रचनाओं~$k_{g,M}$ में भेजता है। अब $f(g) = h_g$ लें। यह
ऐसा फलन है जो $(!A \land !B) \lif !C$ की रचनाओं~$g$ को
$!A \lif (!B \lif !C)$ की रचनाओं में भेजता है। उफ़!{}
\end{ex}

$!A \lor \lnot !A$ कथन को \emph{बहिष्कृत मध्य का नियम} कहते हैं। इसे हम
कुछ विशिष्ट $!A$ के लिए सिद्ध कर सकते हैं (जैसे
$\lfalse \lor \lnot \lfalse$), किंतु सामान्यतः नहीं। कारण यह है कि
अंतःप्रज्ञावादी वियोजन के लिए किसी एक वियोज्य की रचना आवश्यक होती है, पर कुछ
कथन ऐसे हैं जिन्हें अभी न सिद्ध किया जा सकता है, न खंडित (मसलन गोल्डबाख
परिकल्पना)। फिर भी बहिष्कृत मध्य के नियम का खंडन भी नहीं किया जा सकता;
अर्थात $\lnot \lnot (!A \lor \lnot !A)$ सत्य है।

\begin{ex}
$\lnot \lnot (!A \lor \lnot !A)$ सिद्ध करने के लिए हमें ऐसा फलन~$f$ चाहिए
जो $\lnot(!A \lor \lnot !A)$, अर्थात
$(!A \lor (!A \lif \lfalse)) \lif \lfalse$, की किसी रचना को~$\lfalse$
की रचना में बदल दे। दूसरे शब्दों में, हमें ऐसा फलन~$f$ चाहिए कि यदि $g$,
$\lnot(!A \lor \lnot !A)$ की रचना है, तो $f(g)$,~$\lfalse$ की रचना हो।

मान लें कि $g$,~$\lnot(!A \lor \lnot !A)$ की रचना है; अर्थात ऐसा फलन जो
$!A \lor \lnot !A$ की किसी रचना को~$\lfalse$ की रचना में बदलता है।
$!A \lor \lnot !A$ की रचना ऐसा युग्म $\tuple{s, M}$ है जिसमें या तो
$s = 1$ और $M$,~$!A$ की रचना है, अथवा $s = 2$ और $M$,~$\lnot !A$ की
रचना है। $h_1$ वह फलन हो जो~$!A$ की किसी रचना~$M_1$ को
$!A \lor \lnot !A$ की रचना में भेजता है: वह $M_1$ को $\tuple{1, M_2}$ में
भेजता है। और $h_2$ वह फलन हो जो~$\lnot !A$ की किसी रचना~$M_2$ को
$!A \lor \lnot !A$ की रचना में भेजता है: वह $M_2$ को $\tuple{2, M_2}$ में
भेजता है।

$k$ को $\comp{h_1}{g}$ लें: यह ऐसा फलन है जो~$!A$ की कोई रचना दिए जाने पर
$\lfalse$ की रचना लौटाता है; अर्थात यह $!A \lif \lfalse$ अथवा $\lnot !A$
की रचना है। अब $l$ को $\comp{h_2}{g}$ लें। यह ऐसा फलन है जो~$\lnot !A$ की
कोई रचना दिए जाने पर~$\lfalse$ की रचना देता है। चूँकि $k$,~$\lnot !A$ की
रचना है, इसलिए $l(k)$,~$\lfalse$ की रचना है।

समग्रतः हमने यह बताया है कि $\lnot(!A \lor \lnot !A)$ की किसी रचना~$g$ को
$\lfalse$ की रचना में कैसे बदला जा सकता है; अर्थात वह फलन~$f$ जो
$\lnot(!A \lor \lnot !A)$ की रचना~$g$ को~$\lfalse$ की रचना $l(k)$ में
भेजता है, $\lnot\lnot(!A \lor \lnot !A)$ की रचना है।
\end{ex}

जैसा आप देख सकते हैं, !!{formula}ों की अंतःप्रज्ञावादी मान्यता दिखाने के लिए
BHK व्याख्या का उपयोग शीघ्र ही बोझिल और उलझाऊ हो जाता है। सौभाग्य से
अंतःप्रज्ञावादी तर्कशास्त्र के लिए बेहतर !!{derivation}-प्रणालियाँ और अधिक
सटीक अर्थवैज्ञानिक व्याख्याएँ उपलब्ध हैं।
\end{document}
